\documentclass[12pt,a4paper]{article}

% ── Packages ──────────────────────────────────────────────────────────────────
\usepackage[utf8]{inputenc}
\usepackage[T1]{fontenc}
\usepackage{lmodern}
\usepackage{microtype}
\usepackage{amsmath,amssymb,amsthm}
\usepackage{mathtools}
\usepackage{booktabs}
\usepackage{tabularx}
\usepackage{multirow}
\usepackage{array}
\usepackage{xcolor}
\usepackage[colorlinks=true,linkcolor=red!70!black,citecolor=green!50!black,urlcolor=blue!80!black]{hyperref}
\usepackage{graphicx}
\usepackage{geometry}
\usepackage{setspace}
\usepackage{caption}
\usepackage{subcaption}
\usepackage[most]{tcolorbox}
\usepackage{natbib}

\geometry{margin=2.5cm}
\setstretch{1.15}
\graphicspath{{../results/figures/}}

% ── Macros ────────────────────────────────────────────────────────────────────
\newcommand{\SSEE}{\textsc{SSEE}}
\newcommand{\Omm}{\Omega_{\rm m}}
\newcommand{\OmDE}{\Omega_{\rm DE}}
\newcommand{\mnu}{\Sigma m_\nu}
\newcommand{\TODO}[1]{\textcolor{red!70!black}{\textbf{[PENDIENTE R3/R4: #1]}}}

% ── Theorem-like environments ─────────────────────────────────────────────────
\newtheorem{theorem}{Theorem}
\newtheorem{proposition}[theorem]{Proposition}

\title{\textbf{SSEE and the Growth Sector Re-examined:\\
No Second Matter Component, and No $S_8$ Tension,\\
on Raw Survey Data}}

\author{Mike Edison Almeida Vallejo \\
\small \textit{Independent Researcher} \\
\small \href{mailto:mike.almeida1721@gmail.com}{mike.almeida1721@gmail.com} \\
\small ORCID: \href{https://orcid.org/0009-0008-2195-7836}{0009-0008-2195-7836}
}

\date{July 2026 \quad (SSEE, Paper~6 of~10)}

\sloppy
% Numeros generados desde los logs de la cadena de procedencia (src/valores_tex.py)
% ACTA-PROCEDENCIA {"commit": "5651fc97a90dc17d52be0317b46142db9ac86fcb", "reproducible_desde_commit": true, "cambios_sin_commitear": [], "script": "src/valores_tex.py", "script_sha256": "cb5ceb2369bdef5e586d471d57d44a11a9fed354be540a96ce34da85e582efa4", "nucleo_sha256": "5cec8af72d88f154116c5a532e42a9b6ff7a5b4c3d86bcc6921eabb246698993", "entradas": {"manuscript/valores.yaml": "5fa01e91dd68164243d664a73264bb6c87431d2ec0bd69db87c62a983a55992d", "results/logs/kids_publicados.json": "6c25c98059ce91038215358b4e7ac9cfb40198f410354937cee4ea00b46654ff", "results/logs/p10_uv_fscreen.json": "8fb68b295c848e0292e5701e7de795ba7c7ba25e3292d776260bb3dea8970c95", "results/logs/s8_desde_b1.json": "6d225f4403fbfbe6934d7edca26aaefbbe19e05ef6c7b7eda842a259e798b1c1"}, "argv": ["src/valores_tex.py"], "fecha": "2026-09-30T19:35:24", "host": "mike-System-Product-Name", "python": "3.12.3", "versiones": {}, "nucleos": {"OMP_NUM_THREADS": null, "MKL_NUM_THREADS": null, "OPENBLAS_NUM_THREADS": null, "OMPI_COMM_WORLD_SIZE": null}}
% GENERADO por src/valores_tex.py desde manuscript/valores.yaml — NO EDITAR A MANO.
% Cada numero sale de un log de la cadena de procedencia (dvc.yaml + acta).
\providecommand{\val}[1]{\csname ssee@val@#1\endcsname}
\expandafter\def\csname ssee@val@H0_glob\endcsname{67.962142}  % results/logs/p10_uv_fscreen.json#H0_glob
\expandafter\def\csname ssee@val@PTE_k1000_lcdm\endcsname{0.010}  % results/logs/kids_publicados.json#PTE.lcdm.PTE
\expandafter\def\csname ssee@val@PTE_k1000_ssee\endcsname{0.012}  % results/logs/kids_publicados.json#PTE.ssee.PTE
\expandafter\def\csname ssee@val@S8_b1_lcdm\endcsname{0.8229}  % results/logs/s8_desde_b1.json#filas.LCDM.S8
\expandafter\def\csname ssee@val@S8_b1_lcdm_sigma\endcsname{0.0059}  % results/logs/s8_desde_b1.json#filas.LCDM.S8_sigma
\expandafter\def\csname ssee@val@S8_b1_ssee\endcsname{0.8261}  % results/logs/s8_desde_b1.json#filas.SSEE.S8
\expandafter\def\csname ssee@val@S8_b1_ssee_sigma\endcsname{0.0054}  % results/logs/s8_desde_b1.json#filas.SSEE.S8_sigma
\expandafter\def\csname ssee@val@S8_legacy_dato\endcsname{0.8265}  % results/logs/kids_publicados.json#kids_legacy.S8
\expandafter\def\csname ssee@val@S8_legacy_dato_sigma\endcsname{0.0176}  % results/logs/kids_publicados.json#kids_legacy.S8_sigma
\expandafter\def\csname ssee@val@chi2_k1000_lcdm\endcsname{262.75}  % results/logs/kids_publicados.json#PTE.lcdm.chi2
\expandafter\def\csname ssee@val@chi2_k1000_publicado\endcsname{260.31}  % results/logs/kids_publicados.json#kids1000_maxpost.chi2
\expandafter\def\csname ssee@val@chi2_k1000_ssee\endcsname{265.44}  % results/logs/kids_publicados.json#PTE.ssee.chi2
\expandafter\def\csname ssee@val@dAIC_k1000\endcsname{-5.31}  % results/logs/kids_publicados.json#seleccion_kids1000.delta_AIC
\expandafter\def\csname ssee@val@dBIC_k1000\endcsname{-19.0}  % results/logs/kids_publicados.json#seleccion_kids1000.delta_BIC
\expandafter\def\csname ssee@val@dchi2_k1000\endcsname{2.69}  % results/logs/kids_publicados.json#seleccion_kids1000.delta_chi2
\expandafter\def\csname ssee@val@fscreen_IR\endcsname{0.067253}  % results/logs/p10_uv_fscreen.json#f_screen_IR
\expandafter\def\csname ssee@val@fscreen_UV_corr\endcsname{0.002268}  % results/logs/p10_uv_fscreen.json#f_screen_UV_corr
\expandafter\def\csname ssee@val@fscreen_full\endcsname{0.069522}  % results/logs/p10_uv_fscreen.json#f_screen_full
\expandafter\def\csname ssee@val@logA_b1_lcdm\endcsname{3.045}  % results/logs/s8_desde_b1.json#filas.LCDM.logA
\expandafter\def\csname ssee@val@logA_b1_ssee\endcsname{3.042}  % results/logs/s8_desde_b1.json#filas.SSEE.logA
\expandafter\def\csname ssee@val@logA_b1_ssee_sigma\endcsname{0.013}  % results/logs/s8_desde_b1.json#filas.SSEE.logA_sigma
\expandafter\def\csname ssee@val@p_dchi2_k1000\endcsname{0.61}  % results/logs/kids_publicados.json#seleccion_kids1000.p_delta_chi2
\expandafter\def\csname ssee@val@sK_IR\endcsname{0.403302}  % results/logs/p10_uv_fscreen.json#s_K_IR
\expandafter\def\csname ssee@val@sK_UV_corr\endcsname{0.013603}  % results/logs/p10_uv_fscreen.json#s_K_UV_corr
\expandafter\def\csname ssee@val@sK_full\endcsname{0.416905}  % results/logs/p10_uv_fscreen.json#s_K_full
\expandafter\def\csname ssee@val@tension_b1_lcdm\endcsname{2.59}  % results/logs/s8_desde_b1.json#filas.LCDM.tension_kids
\expandafter\def\csname ssee@val@tension_b1_ssee\endcsname{2.73}  % results/logs/s8_desde_b1.json#filas.SSEE.tension_kids
\expandafter\def\csname ssee@val@z_k1000_lcdm\endcsname{2.5}  % results/logs/kids_publicados.json#PTE.lcdm.z
\expandafter\def\csname ssee@val@z_k1000_ssee\endcsname{2.4}  % results/logs/kids_publicados.json#PTE.ssee.z

\begin{document}
\maketitle

\begin{abstract}
We test the growth-of-structure sector of \SSEE{} against \emph{raw} survey
data rather than against compressed statistics whose reduction assumes a
$\Lambda$CDM background. The model has a single matter density,
$\Omega_m=\val{OMEGA_M_TOTAL_d6}$, entering the background and the growth of structure
alike; the background is fixed by the algebra. Against the 357 raw
$\xi_\pm$ points of KiDS-Legacy \citep{Wright2025}, with the amplitude free, the
shear prefers $\log(10^{10}A_s)=3.0255\pm0.0396$, $0.49\sigma$ from the
$3.04483$ that \SSEE{}'s own CMB fit supplies (under the identical pipeline
with the Planck background, $1.01\sigma$). Fixing the amplitude to the CMB
value leaves the growth sector with \emph{no free cosmological parameter} and
costs nothing: $\chi^2_{\rm min}=417.97$ against $418.34$ with the amplitude
free, $\Delta\mathrm{BIC}=\val{dBIC_leg_libre}$, and the nuisance parameters do not move to
compensate. $S_8=\val{S8_leg_pred}$ is then an output of the algebra and the CMB, against
a measured $\val{S8_legacy_dato}\pm\val{S8_legacy_dato_sigma}$. KiDS-Legacy predates this analysis by sixteen
months and no temporal priority is claimed: the result is a consistency of
$A_s$ between epochs. We state against our own interest that the three
$\chi^2$ values compared span less than $1.0$: cosmic shear does not
discriminate between these backgrounds, and what separates them is the
parameter count. On the earlier KiDS-1000 release (225 raw points, amplitude
free), \SSEE{} gives $\chi^2_{\rm min}=\val{chi2_k1000_ssee}$ against $\val{chi2_k1000_lcdm}$ for
$\Lambda$CDM run through the identical pipeline with its background free: four
extra parameters buy $\Delta\chi^2=\val{dchi2_k1000}$, less than the $4$ they would return
by chance ($\Delta\mathrm{BIC}=\val{dBIC_k1000}$ in favour of \SSEE{}). Against the raw
BOSS DR12 multipoles (222 points, one-loop LPT) the algebraic and Planck
backgrounds fit equally well within noise, $\Delta\chi^2=\val{bossR_dchi2}$ for \SSEE{} at equal freedom; clustering does not
constrain the amplitude. This paper also retracts a second matter sector and
an associated particle ($m_\phi=40.70$~eV) proposed in an earlier version: the
sector was defined by subtracting $1+w_0\approx0.160$ --- a number of the
dark-energy equation of state, not a density --- from the matter density, and
the $S_8$ tension it was built to close existed only in the compressed
statistic, never in the raw data.
\end{abstract}

\clearpage
{\small\tableofcontents}
\clearpage

% ============================================================================
\section{Introduction: why the growth sector needed re-examination}
% ============================================================================

The background sector of \SSEE{} --- the algebraic derivation of $H_0$,
$\omega_b$, $\omega_c$, $n_s$, $w_0$ and $w_a$ from $\varphi$ and $\pi$ --- has
been tested against Planck PR4, DESI DR2 BAO, and Pantheon+ in Papers~1--4,
9--10. The growth-of-structure sector --- how matter perturbations evolve once
the background is fixed --- is tested independently, against two late-time
probes: weak gravitational lensing (cosmic shear, KiDS-1000) and redshift-space
galaxy clustering (BOSS DR12). Both probes report their results as derived
quantities: KiDS-1000 publishes $S_8\equiv\sigma_8\sqrt{\Omega_m/0.3}$, and BOSS
publishes $f\sigma_8(z)$, the growth rate times the amplitude of clustering.

Both derived quantities share a hidden dependency that matters for a model
whose background departs from $\Lambda$CDM: converting a raw two-point
statistic (shear correlation functions, or galaxy power-spectrum multipoles)
into a compressed number requires assuming a fiducial cosmology --- the
Alcock--Paczynski effect \citep{AlcockPaczynski1979} for redshift-space
clustering, and the non-linear matter power spectrum template for lensing.
Both official pipelines assumed $\Lambda$CDM calibrated to Planck
\citep{Planck2020}. Comparing a $\Lambda$CDM-conditioned number against a
background that is \emph{not} $\Lambda$CDM (\SSEE{} has a
Chevallier--Polarski--Linder \citep{ChevallierPolarski2001,Linder2003}
equation of state $w_0=\val{W0_d6}$, not $-1$) imports part of the $\Lambda$CDM
background into the comparison itself.

This paper re-does both measurements on the raw data products --- KiDS-1000
$\xi_\pm(\theta)$ (Asgari et al.~2021 \citep{Asgari2021}) and BOSS DR12 power-spectrum multipoles
$P_0,P_2,P_4(k)$ (Beutler et al.~2017 \citep{Beutler2017}) --- applying the model-dependent
deformation (Alcock--Paczynski for BOSS; the non-linear matter power spectrum
template for KiDS) to the \emph{theory prediction}, never to the data. This
keeps the measurement clean of any assumed background, at the cost of
re-deriving both likelihoods from the covariance matrices up rather than
reusing the compressed public numbers.

% ============================================================================
\section{The retraction: why the second sector is withdrawn}
\label{sec:retraction}
% ============================================================================

\subsection{The subtraction that never was a density}

The two-sector construction, retired here, rested on
\begin{equation}
  \Omega_{\phi\rm DM} \;=\; \Omega_{m,\rm CMB} - \Omega_{m,\rm dyn},
  \qquad \Omega_{m,\rm dyn}\equiv 1+w_0.
\end{equation}
$\Omega_{m,\rm CMB}=\omega_m/h^2=\val{OMEGA_M_TOTAL_d6}$ is a measured density fraction, fixed
by the CMB-calibrated background. $\Omega_{m,\rm dyn}=1+w_0\approx0.160050$ is,
by construction, minus the dark-energy equation of state --- an algebraic
identity of the background sector's $w_0=-T_R/M_V$, entering the Friedmann
equations as an equation-of-state parameter, never as a density. Both
quantities are dimensionless, so the subtraction is arithmetically well
defined; but a well-defined subtraction between a density and an
equation-of-state number does not itself define a density, let alone a
particle. The internal code comments that introduced this construction stated
explicitly that $\Omega_{m,\rm dyn}$ ``is NOT a matter density: it is the
dynamical cold sector ($1+w_0$)'' in the same breath as summing it
with $\Omega_{\phi\rm DM}$ to recover $\Omega_{m,\rm CMB}$ --- an internal
inconsistency once both statements are read together. ($1+w_0$ is
dimensionless and is, by construction, minus the dark-energy equation of
state; the parenthetical gloss is ours, not the source comment's.)

\subsection{Falsifiability is not evidence of existence}

A canonical particle mass, $m_\phi=40.70$~eV --- retracted here --- was
derived from this subtraction via a fixed algebraic multiplier
(retracted: $\rm SOLAR^2\!\cdot\!KRYSTOS_V=594.28$) acting on $\Sigma m_\nu^{\rm active}$,
explicitly to render the two-sector hypothesis falsifiable against
free-streaming suppression in the matter power spectrum. Falsifiability is a
necessary property of a scientific hypothesis; it is not evidence that the
hypothesis's object exists. Once the subtraction that defined
$\Omega_{\phi\rm DM}$ is shown not to correspond to a physical density, the
particle built on top of it has nothing left to be made of, and is withdrawn
together with it --- independently of whatever bound the data would place on
its mass. \S\ref{sec:s8} reports that bound only as part of the historical
record of how the retraction was reached, not as a surviving result.

\subsection{Open problems closed by dissolution}

Four open problems catalogued against the retired two-sector construction close by
dissolution, following the precedent of OP-8 (the analogous $\Omega_m$
partition via the algebraic constant MIRA, retired 2026-07-09):

\begin{itemize}
  \item \textbf{OP-9} (origin of the SOLAR$^2\cdot$KRYSTOS$_V$ coefficient) ---
        closes: there is no coefficient left to derive.
  \item \textbf{OP-10} (the ``single free scalar field'' claim, admitted as
        ``technically violated'' by the retired two-sector construction) --- closes,
        and the single-field philosophy of \SSEE{} is restored without
        qualification.
  \item \textbf{OP-11} ($\xi$, the non-minimal coupling of the second sector
        --- an explicit additional free continuous parameter) --- closes: the
        parameter is removed, not fitted.
  \item \textbf{OP-12} ($T_\phi$, described in its own documentation as
        ``bookkeeping, not a temperature derived from physics'') --- closes.
\end{itemize}

No new open problem is created by this retraction: what remains is
$\Omega_m=\val{OMEGA_M_TOTAL_d6}$, a single measured quantity, entering both the background
and the growth-of-structure sector with no partition.

% ============================================================================
\section{Methodology: measuring against raw data, not compressed statistics}
\label{sec:method}
% ============================================================================

\subsection{The principle: combined data reports its own internal tensions}

A dataset built by combining independent probes under a flexible model does
not report fidelity to any one of them: it reports the point at which their
mutual tensions are minimized, which is not the same thing. A model with
enough freedom to fit each probe separately with a \emph{different} set of
nuisance values will inevitably report a lower $\chi^2$ when its parameters are
free to migrate between probes than when they are held fixed to what any
single probe prefers ---not because the combination is more accurate, but
because it is averaging disagreement rather than resolving it. This is the
generic mechanism behind the well-documented $S_8$ tension between Planck and
cosmic shear surveys under $\Lambda$CDM \citep{DiValentino2021}: parameters
fitted to the CMB \citep{Planck2020}, carried over unchanged to lensing, land
outside the region preferred by KiDS-1000 \citep{Asgari2021,Heymans2021} and
DES Y3 \citep{Amon2022} alike.

\SSEE{} is structurally unable to exploit this mechanism: its background
(\(\omega_b,\omega_c,h,n_s,w_0,w_a\)) is fixed by algebra, identical whether
evaluated against the CMB, cosmic shear, or redshift-space clustering. The
only cosmological parameter that could vary between probes is the
primordial amplitude $A_s$ (equivalently $\log A$). In \SSEE{} it is counted
once, at the CMB, and carried to every other probe fixed
(\S\ref{sec:legacy}).

\subsection{Test design: ingredients transported, not re-fitted}
\label{sec:r7}

To make this principle falsifiable rather than rhetorical, we perform the
following test (a design due to M.~Almeida): the three probes --- Planck
CMB, KiDS-1000 shear, BOSS DR12 clustering --- are combined with the shared
amplitude free, yielding the ingredient set the \emph{combination} prefers.
That same ingredient set is then evaluated against each probe
\emph{individually}, without re-fitting. Transporting parameters from a joint
fit to an individual dataset always degrades $\chi^2$ relative to a dedicated
fit --- that is arithmetic, true of any model, and proves nothing by itself.
The diagnostic quantity is how much it degrades relative to what the
probe's own statistical uncertainty predicts:
\begin{equation}
  \Delta\chi^2(\text{probe}) \;=\;
  \chi^2\big(\text{joint ingredients}\big) - \chi^2\big(\text{probe-fitted ingredients}\big).
\end{equation}
For \SSEE{}, $\Delta\chi^2\equiv0$ by construction on every background
parameter except $A_s$: there is no ingredient to transport, because none of
$\omega_b,\omega_c,h,n_s,w_0,w_a$ are free to begin with. For $\Lambda$CDM,
$\Delta\chi^2>0$ is measured, not assumed. Results in
Table~\ref{tab:r7}.

\begin{table}[h]
\centering
\caption{Ingredient transport test (\S\ref{sec:r7}). $\Delta\chi^2$ evaluated
per probe, joint-fit ingredients vs.\ probe-dedicated fit.}
\label{tab:r7}
\begin{tabular}{lccc}
\toprule
 & CMB (Planck) & shear (KiDS-1000) & clustering (BOSS) \\
\midrule
$\Lambda$CDM, $\Delta\chi^2$ & \TODO{joint fit} & \TODO{joint fit} & \TODO{joint fit + free-background BOSS} \\
\SSEE{}, $\Delta\chi^2$ & 0 (by construction) & 0 (by construction) & 0 (by construction) \\
\bottomrule
\end{tabular}
\end{table}

\subsection{Raw KiDS-1000 cosmic shear likelihood}
\label{sec:kids-method}

We evaluate $\chi^2$ against the 225 (of 270) scale-cut $\xi_\pm(\theta)$ data
points from the KiDS-1000 Blind~C release (Asgari et al.~2021 \citep{Asgari2021}), using the flat-sky
Limber approximation for the shear power spectrum $C_\ell^{ij}$ (GG+GI+II, NLA
intrinsic alignment model, Bridle \& King 2007 \citep{BridleKing2007}), HMcode-2015 non-linear
matter power spectrum (Mead et al.~2015 \citep{Mead2015}) with the one-parameter baryon
feedback prescription used in the official KiDS-1000 analysis, and the
official redshift-distribution SOM covariance and $\delta_c$ priors. The
pipeline reproduces the published maximum-posterior likelihood to 0.4\%
(negative control: $\chi^2=261.27$ vs.\ \val{chi2_k1000_publicado} official, the latter from the
public nested-sampling chain, MultiNest \citep{Feroz2009}).

Eight nuisance parameters (baryon feedback amplitude, intrinsic alignment
amplitude, five source-redshift shifts, additive $c$-term) are marginalized
via MCMC (Cobaya \citep{Cobaya}, Metropolis sampler, 4 independent chains,
Gelman--Rubin $R-1<0.03$), not minimized, for both \SSEE{} (background fixed,
$\log A$ free: 9 parameters total) and $\Lambda$CDM (background free:
$\omega_b,\omega_c,h,n_s$; 13 parameters total, same nuisance treatment).

\subsection{Raw BOSS DR12 redshift-space multipoles}
\label{sec:boss-method}

We evaluate $\chi^2$ against the full-shape power-spectrum multipoles
$P_0,P_2,P_4(k)$ from BOSS DR12 (Beutler et al.~2017 \citep{Beutler2017}), 3 redshift bins
($z=0.38,0.51,0.61$) $\times$ 2 hemispheres (NGC/SGC), using the PATCHY mock
covariance with Hartlap correction. The Alcock--Paczynski deformation is
applied to the \emph{model}, mapped onto the fiducial cosmology
($\Omega_m=0.31$, $w=-1$, $H_0=100$, declared in the data headers) in which
Beutler measured, so that the raw data vector is never touched. The survey
window function is applied via the configuration-space convolution method of
Wilson et al.~(2017) \citep{Wilson2017}: $P_\ell(k)\to\xi_\ell(s)$ (Hankel transform),
mixed with the window multipoles $W_0,W_2,W_4,W_6,W_8$, transformed back.
Redshift-space distortions are modelled at one loop in Lagrangian
perturbation theory, in its effective-field-theory implementation
(velocileptors; \citealp{Baumann2012,ChenVlahWhite2020}), to
$k_{\max}=0.20\,h\,$Mpc$^{-1}$; an earlier linear-Kaiser pass is described in
\S\ref{sec:limitation-kaiser} and does not publish.

% ============================================================================
\section{$S_8$ against raw cosmic shear}
\label{sec:s8}
% ============================================================================

The \SSEE{} chain (R3) samples nine free parameters --- the primordial
amplitude $\log A$ plus the eight KiDS-1000 nuisance parameters (the baryonic
feedback amplitude $A_{\rm baryon}$, the intrinsic-alignment amplitude
$A_{\rm IA}$, five photometric-redshift shifts $\delta z_i$, and the
additive shear term $\delta_c$) --- against the 225 unmasked
$\xi_\pm$ data points, with the background held fixed by the algebra
($\Omega_m=\val{OMEGA_M_TOTAL_d6}$, $\omega_b$, $\omega_c$, $h$, $n_s$, $w_0$, $w_a$ all
determined, none of them fitted). Four independent MPI chains were run to a
Gelman--Rubin convergence of $R-1=\val{Rm1_k1000}$, well inside the stopping criterion
$R-1<0.03$, giving $\val{neff_k1000}$ effective samples after discarding the first 30\% of each chain.
The marginalized posterior is
\begin{equation}
\sigma_8 = \val{sigma8_k1000_ssee}\pm\val{sigma8_k1000_ssee_sigma}, \qquad
S_8 \equiv \sigma_8\sqrt{\Omega_m/0.3} = \val{S8_k1000_ssee}\pm\val{S8_k1000_ssee_sigma},
\end{equation}
with a best-fit $\chi^2=\val{chi2_k1000_ssee}$ over 225 data points and 9 free parameters
(216 degrees of freedom, $\chi^2/{\rm dof}=1.23$). Both nuisance parameters
that carry astrophysical meaning land in unremarkable places
($A_{\rm baryon}=2.57\pm0.31$, $A_{\rm IA}=0.55\pm0.29$), so the fit is not
being rescued by a nuisance parameter running to the edge of its prior.

\paragraph{Where the model is actually tested: the fit to the raw data.}
$S_8$ is an output of a fit, not a measurement: quoting one model's $S_8$
against another's compares two \emph{products} of two analyses. The
quantity that confronts \SSEE{} with something the model did not produce is
the likelihood of the 225 raw $\xi_\pm$ points, so that is where we put the
comparison (Table~\ref{tab:s8}). Against the $\Lambda$CDM control run
through the same pipeline, same mask, same covariance and same nuisance
treatment,
\begin{equation}
\Delta\chi^2 \;\equiv\; \chi^2_{\SSEE{}}-\chi^2_{\Lambda\rm CDM}
\;=\; \val{chi2_k1000_ssee}-\val{chi2_k1000_lcdm} \;=\; \val{dchi2_k1000} ,
\qquad \Delta k = 4 ,
\label{eq:dchi2}
\end{equation}
where $\Delta k$ counts the parameters $\Lambda$CDM has and \SSEE{} does not.
Four additional free parameters that carried no information would still
lower $\chi^2$ by $\langle\Delta\chi^2\rangle=\Delta k=4$ on average; here
they lower it by $\val{dchi2_k1000}$, and $p(\Delta\chi^2\ge \val{dchi2_k1000}\,|\,\Delta k=4)=\val{p_dchi2_k1000}$.
The extra background freedom of $\Lambda$CDM is therefore not merely
unnecessary on this dataset --- it is statistically indistinguishable from
absent. Penalising it explicitly gives $\Delta\mathrm{AIC}=\val{dAIC_k1000}$ and
$\Delta\mathrm{BIC}=\val{dBIC_k1000}$, both favouring \SSEE{}, the latter in the
\emph{strong} range on the Jeffreys scale.

The derived amplitudes are then a consequence rather than the test:
$\Delta S_8 = 0.0016$, under a tenth of either error bar. Against the
\emph{published} KiDS-1000 value $S_8=\val{S8_k1000_dato}\pm\val{S8_k1000_dato_sigma}$ the offset is
$\val{S8_k1000_ssee_tension}\sigma$, but that published number is itself $\Lambda$CDM-derived, so
this second comparison is informative about the pipeline rather than about
the cosmology, and we do not rest any claim on it.

\paragraph{The one tension we do report, and it is shared.} Both fits are
formally poor in absolute terms: $\chi^2=\val{chi2_k1000_ssee}$ against an expectation of
$216\pm20.8$ for \SSEE{} and $\val{chi2_k1000_lcdm}$ against $212\pm20.6$ for
$\Lambda$CDM, i.e.\ $\val{z_k1000_ssee}\sigma$ and $\val{z_k1000_lcdm}\sigma$ high, with probabilities to
exceed of $\val{PTE_k1000_ssee}$ and $\val{PTE_k1000_lcdm}$. This is not a feature of either cosmology:
the KiDS-1000 published best fit on these same points has
$\chi^2=\val{chi2_k1000_publicado}$, which under any of the plausible effective-parameter counts
gives a comparable $p\sim0.01$--$0.03$. The excess is a property of the
$\xi_\pm$ data vector --- plausibly residual small-scale modelling or
covariance calibration --- and it is inherited identically by every model
fitted to it. We flag it rather than absorb it, and note that our control's
reproduction of the official value ($\val{chi2_k1000_lcdm}$ vs.\ $\val{chi2_k1000_publicado}$, $0.9\%$) is what
demonstrates the excess is not introduced by our implementation.

\paragraph{What this does and does not claim.} It would be wrong to read this
as \SSEE{} having carried a $3.5\sigma$ tension and then removing it.
\SSEE{} never carried that tension. At this level the model has
exactly \emph{one} free quantity, the primordial amplitude $A_s$; the
background is fixed by the algebra and is identical in every fit reported here.
What varies between determinations is only \emph{which dataset is asked to fix
$A_s$}, and the datasets do not agree with each other:

\begin{center}
\begin{tabular}{lccc}
\toprule
$A_s$ determined by & $\log(10^{10}A_s)$ & $S_8$ & vs.\ published KiDS-1000$^\dagger$ \\
\midrule
raw KiDS-1000 $\xi_\pm$ (this work) & $\val{logA_k1000}\pm\val{logA_k1000_sigma}$ & $\val{S8_k1000_ssee}\pm\val{S8_k1000_ssee_sigma}$ & $\val{S8_k1000_ssee_tension}\sigma$ \\
\SSEE{} fit to raw Planck (full \texttt{plik}, B1 chains) & $\val{logA_b1_ssee}\pm\val{logA_b1_ssee_sigma}$ & $\val{S8_b1_ssee}\pm\val{S8_b1_ssee_sigma}$ & $\val{tension_b1_ssee}\sigma$ \\
$\Lambda$CDM-fitted Planck value (borrowed) & $\val{logA_b1_lcdm}$ & $\val{S8_b1_lcdm}\pm\val{S8_b1_lcdm_sigma}$ & $\val{tension_b1_lcdm}\sigma$ \\
\bottomrule
\end{tabular}

\smallskip
\footnotesize
$^\dagger$The last column compares each row against the \emph{published}
KiDS-1000 value, which is a $\Lambda$CDM-derived number; it is given for
orientation only and no claim here rests on it. The quantity that matters for
the argument below is different: the internal disagreement between rows~1
and~2, which are the same model with the same background. In
$\log(10^{10}A_s)$ that separation is $\val{logA_b1_ssee}-\val{logA_k1000}=\val{logA_k1000_vs_b1_dif}\pm\val{logA_k1000_vs_b1_dif_sigma}$, i.e.\
$\mathbf{\val{logA_k1000_vs_b1}\sigma}$.
\normalsize
\end{center}

The first two rows are the same model with the same background, differing only
in the data used to set one number, and they sit $\val{logA_k1000_vs_b1}\sigma$ apart. That
separation is a disagreement \emph{between the datasets}, not a property of
\SSEE{}: it persists when $A_s$ is fixed by \SSEE{}'s \emph{own} CMB fit rather
than borrowed from $\Lambda$CDM, and it is the same Planck--shear discrepancy
that $\Lambda$CDM also carries. \textbf{We do not claim to resolve it.}

What the earlier version of this paper did wrong was to attribute that
data--data disagreement to the model --- reporting it as ``\SSEE{}'s $S_8$
tension'' --- and then to build a second matter sector and a particle to close
it. The disagreement was never the model's to close, and the KiDS
collaboration's own recalibration dissolved it in the data (\S\ref{sec:legacy}).

% ============================================================================
\section{The same test against KiDS-Legacy}
\label{sec:legacy}
% ============================================================================

The disagreement reported above was resolved by the KiDS collaboration itself,
and not in the model's favour or against it: in the data. KiDS-Legacy
\citep{Wright2025} recalibrates the source redshift distributions and re-analyses
the complete survey, and its Appendix~I re-runs the KiDS-1000 analysis end to
end with the new methods under blinding. The tension against Planck dissolves in
steps, quoted there as signed consistencies in $S_8$ (negative meaning KiDS
below Planck): $-2.39\sigma \to -1.42\sigma$ from the $n(z)$ calibration
alone, then $-1.12\sigma$, $-0.85\sigma$ and $-0.74\sigma$ as tomography,
alignment model and sky area are updated. Their stated threshold is
``the typical KiDS threshold for consistency of $2.3\sigma$ (equivalent to a
one-sided probability to exceed of $0.01$)''; every step with the new
calibration falls below it.

\paragraph{Chronology, stated plainly.} KiDS-Legacy appeared on arXiv on
2025 March 25, sixteen months before the analysis of \S\ref{sec:s8} was written.
\textbf{We claim no temporal priority over it.} What we report here is an
internal consistency test of this model, run on the raw Legacy data vector, and
the reader should weigh it as such.

We repeat the measurement of \S\ref{sec:s8} on the 357 unmasked $\xi_\pm$ points
of the Legacy release, in three configurations that differ only in what is held
fixed. All three use the same likelihood, the same scale cuts
($\xi_+$: $2'$--$300'$, $\xi_-$: $4'$--$300'$) and the same eight nuisance
parameters.

\begin{center}
\footnotesize
\begin{tabular}{lccccc}
\toprule
Configuration & background & $A_s$ & free & $\chi^2_{\rm min}$ & BIC \\
\midrule
\SSEE{} unified & algebraic & from CMB & 8 & $\val{chi2_leg_unif}$ & $\val{BIC_leg_unif}$ \\
\SSEE{}         & algebraic & free     & 9 & $\val{chi2_leg_libre}$ & $\val{BIC_leg_libre}$ \\
$\Lambda$CDM-fixed & Planck 2018 & free & 9 & $\val{chi2_leg_lcdm}$ & $\val{BIC_leg_lcdm}$ \\
\bottomrule
\end{tabular}
\normalsize
\end{center}

Three results follow, and they answer different questions.

\paragraph{(i) The amplitude the shear asks for is the one the CMB fixes.}
With the background held by the algebra and $\log A$ free, the Legacy data
prefer $\log(10^{10}A_s) = 3.0255\pm0.0396$, against the value \SSEE{}'s own CMB
fit supplies, $3.04483$: a separation of $0.49\sigma$. Under the identical
pipeline with the Planck background held fixed instead, the same data prefer
$3.0069\pm0.0368$ against that background's own CMB value, $1.01\sigma$ away.
The difference has a mechanism rather than being noise: \SSEE{}'s $\Omega_m$ is
$2.00\%$ lower than Planck's, and the shear asks it for an $A_s$ $1.88\%$
higher --- less matter requires more amplitude to produce the same observed
clustering. Each background \emph{forces} a different $A_s$, and the one the
algebraic background forces is the one its own CMB predicts. Both $\log A$
values come from the same 357 points and are therefore correlated; no clean
significance can be quoted for their difference, and $1.01\sigma$ is not a
tension either. What can be stated is the direction, and that it has an
identified physical cause.

\paragraph{(ii) Fixing $A_s$ costs nothing, and could have cost a great deal.}
Removing the only cosmological parameter from a fit removes its ability to
accommodate: if either the algebraic background or the CMB amplitude were wrong,
$\chi^2$ rises with nowhere to hide it. It did not rise. With one parameter
fewer, $\chi^2_{\rm min}=417.97$ against $418.34$, and
$\Delta\mathrm{BIC}=\val{dBIC_leg_libre}$ in favour of the unified configuration
($\val{dBIC_leg_lcdm}$ against $\Lambda$CDM-fixed). A second, independent indication points
the same way: the nuisance parameters do not move to compensate. Had the fixed
$A_s$ been wrong, $\log T_{\rm AGN}$ and $A_{\rm scale}$ would have run to cover
the gap; they shift by $0.066$ and $0.10$, well inside their own posterior
widths. We note for completeness that $\Delta\chi^2$ is formally negative
($-0.365$), which cannot be a true improvement --- the fixed-parameter space is
contained in the free one --- and measures the MCMC sampling noise: eight
dimensions are explored more efficiently than nine. The honest statement is
that the cost is indistinguishable from zero, not that fixing improves the fit.

\paragraph{(iii) $S_8$ becomes a prediction.} With both background and amplitude
supplied, $S_8=\val{S8_leg_pred}$ is an output of the algebra and the CMB, with no free
cosmological parameter behind it and therefore no posterior width of its own.
The official Legacy $\xi_\pm$ posterior, computed from its own published chain
as $\sigma_8\sqrt{\Omega_m/0.3}$, gives $\val{S8_legacy_dato}\pm\val{S8_legacy_dato_sigma}$. Two caveats bound
what this establishes: the agreement in $S_8$ and the $0.49\sigma$ in $\log A$
are the same fact in two variables and do not compound, and the $\pm0.0176$
width means any prediction between $0.809$ and $0.844$ would have agreed. The
test is not sharp. What it does establish is that the growth sector takes its
amplitude from the CMB and describes cosmic shear without adjusting anything.

\paragraph{What this closes.} The gap that the retracted light particle was
introduced to fill is $0.49\sigma$ wide. There is no work left for it to do, and
it is withdrawn as unnecessary --- not excluded by a bound. Two things follow
that are worth stating against our own interest. First, a particle that is not
needed may still exist; this is not a constraint on its mass. Second, the three
$\chi^2$ values above span less than $1.0$, so cosmic shear does not discriminate
between these backgrounds: what separates them is the parameter count (BIC) and
the cross-epoch consistency of $A_s$, neither of which is a goodness-of-fit
statement. \textbf{The result is not that \SSEE{} describes cosmic shear better;
it is that \SSEE{} describes it equally well without being allowed to fit it.}

\begin{table}[h]
\centering
\caption{Fit to the \emph{raw} KiDS-1000 $\xi_\pm$ data vector (225 points),
primordial amplitude free in both models, MCMC-marginalized. The test is the
left block: the likelihood of data neither model produced. The $S_8$ block is
derived output, shown for continuity with the literature. $p$ is the
probability to exceed of $\chi^2_{\rm min}$ at the stated dof.}
\label{tab:s8}
\begin{tabular}{lccccccc}
\toprule
 & $k$ & dof & $\chi^2_{\rm min}$ & $\chi^2/$dof & $p$ & $\sigma_8$ & $S_8$ \\
\midrule
$\Lambda$CDM (background free)  & 13 & 212 & $\val{chi2_k1000_lcdm}$ & $1.239$ & $\val{PTE_k1000_lcdm}$ & $0.875\pm0.089$ & $\val{S8_k1000_lcdm}\pm\val{S8_k1000_lcdm_sigma}$ \\
\SSEE{} (background fixed)      &  9 & 216 & $\val{chi2_k1000_ssee}$ & $1.229$ & $\val{PTE_k1000_ssee}$ & $\val{sigma8_k1000_ssee}\pm\val{sigma8_k1000_ssee_sigma}$ & $\val{S8_k1000_ssee}\pm\val{S8_k1000_ssee_sigma}$ \\
\midrule
\multicolumn{8}{l}{\footnotesize $\Delta\chi^2=\val{dchi2_k1000}$ for $\Delta k=4$ (expected $4$ if the extra parameters carry no information, $p=\val{p_dchi2_k1000}$);} \\
\multicolumn{8}{l}{\footnotesize $\Delta\mathrm{AIC}=\val{dAIC_k1000}$, $\Delta\mathrm{BIC}=\val{dBIC_k1000}$, both favouring \SSEE{}.} \\
\midrule
KiDS-1000 published & --- & --- & $\val{chi2_k1000_publicado}$ & --- & --- & --- & $\val{S8_k1000_dato}\pm\val{S8_k1000_dato_sigma}$ \\
\bottomrule
\end{tabular}
\end{table}

The $\Lambda$CDM row is the methodological control (R4): the same pipeline,
the same raw data, the same nuisance treatment, with the four background
parameters $\omega_b,\omega_c,h,n_s$ additionally free. Its purpose is not to
provide a new measurement --- KiDS-1000 has already published theirs --- but
to demonstrate that this pipeline reproduces the published $\Lambda$CDM
result, which is what licenses the \SSEE{} row above it. It does:
$S_8=\val{S8_k1000_lcdm}\pm\val{S8_k1000_lcdm_sigma}$ against the published $\val{S8_k1000_dato}\pm\val{S8_k1000_dato_sigma}$. We stress that
this agreement is \emph{not} a physical measurement of consistency. The
published KiDS-1000 value is itself a $\Lambda$CDM number --- it is what one
obtains by fitting $\Lambda$CDM to these data --- so a $\Lambda$CDM run
reproducing it tests the code path, not the cosmology. Reading it as a
``tension'' would be a category error.

What the control licenses is the comparison in Table~\ref{tab:s8} between the
two rows above the rule, which is like for like: identical data, identical
code, identical nuisance treatment. The comparison that carries the weight is
the $\chi^2$ block, Eq.~\eqref{eq:dchi2}: \SSEE{} pays $\val{dchi2_k1000}$ in $\chi^2$ for
giving up four parameters, where four uninformative parameters would have been
worth $4$. The $\Delta S_8 = 0.0016$ in the right-hand block is a restatement
of the same fact in derived quantities --- under a tenth of either error bar,
and, since both fits are to the \emph{same} data, with strongly correlated
uncertainties that make the true discrepancy smaller still. \SSEE{} therefore
describes the raw shear data as well as $\Lambda$CDM does, using four fewer
free parameters and with its background fixed algebraically rather than
fitted.

Two entries in the $\Lambda$CDM row must not be quoted in isolation:
$\sigma_8=0.875\pm0.089$ and $\Omega_m=0.231\pm0.046$ (and correspondingly the
broad $\log A$) are individually poorly determined, because cosmic shear
constrains essentially only their combination $S_8$ --- which is why that
combination is the reported statistic. The published KiDS-1000 analysis shows
the same degeneracy.

\subsection{Historical note: the particle's exclusion, for the record}

Before the subtraction defining $\Omega_{\phi\rm DM}$ was recognized as
physically empty (\S\ref{sec:retraction}), the two-sector hypothesis was
tested directly against raw KiDS-1000 data with the now-retracted particle mass
$m_\phi=40.70$~eV. Free-streaming suppresses $\sigma_8$ by 10.4\% at that
mass, against a raw-data uncertainty of 3.16\%; the amplitude $A_s$ cannot
compensate because free-streaming changes the \emph{shape} of $P(k)$, not its
overall scale. A measured suppression law,%
% ORIGEN-VALOR: 2.283 — exponente de la ley de supresion de la PARTICULA RETIRADA (julio 2026); el script que la ajusto no se guardo. Parrafo historico: se cita como registro de la retractacion, no como resultado vigente
$\Delta\sigma_8(\%)=52088\cdot m_\phi^{-2.283}$, places a lower bound
$m_\phi>70.3$~eV for consistency with the shear data --- 1.73$\times$ the
canonical mass, and outside the maximal mass reachable by the model's own
algebraic lineage grammar (§\ref{sec:retraction}, maximum multiplier 594.28).
We record this bound as part of how the retraction was reached
(\S\ref{sec:retraction}), not as a live result: the particle it would
constrain does not exist in the current model.

% ============================================================================
\section{$f\sigma_8$: growth rate against raw BOSS DR12 clustering}
\label{sec:fsigma8}
% ============================================================================

The clustering counterpart of \S\ref{sec:s8} asks the same question of the
other growth probe: with the background held fixed --- algebraically in
\SSEE{}, at Planck 2018 in $\Lambda$CDM --- what amplitude does the raw
galaxy clustering prefer, and at what $\chi^2$? The test is the $\chi^2$ over
the 222 multipole points, \emph{not} $f\sigma_8$: $f\sigma_8$ is an output of
the fit, and therefore a product of the model, the same lesson $S_8$ taught in
\S\ref{sec:retraction}.

Both models carry identical freedom here: one shared amplitude
$\ln(10^{10}A_s)$ plus $(b_1,b_2,b_s,\alpha_0,\alpha_2,{\rm SN_0})$ per
$(z\text{-bin},\text{hemisphere})$, $37$ in total, of which the $18$ that
enter the model linearly are marginalised in closed form
(\S\ref{sec:boss-method}) and $19$ are sampled. So $\Delta k=0$: this
comparison carries no parsimony credit in either direction. \SSEE{}'s
advantage in parameter count lives at the CMB ($k=2$ against $6$) and in the
shear likelihood ($9$ against $13$), not here. What this run answers is
narrower and cleaner: \emph{with the same nuisance freedom, does the algebraic
background fit raw clustering as well as Planck's?}

\begin{table}[h]
\centering
\caption{$f\sigma_8$ against raw BOSS DR12 multipoles (R1/R2), LPT to
$k_{\max}=0.20\,h\,$Mpc$^{-1}$. Both backgrounds fixed; $222$ data points,
$37$ free parameters, $185$ degrees of freedom. Published values are the
BOSS DR12 consensus of \citet{Alam2017}.}
\label{tab:fsigma8}
\begin{tabular}{lccc}
\toprule
 & $z=0.38$ & $z=0.51$ & $z=0.61$ \\
\midrule
\SSEE{}                    & $\val{bossR_fs8_ssee_z1}\pm\val{bossR_fs8_ssee_z1_sigma}$ & $\val{bossR_fs8_ssee_z2}\pm\val{bossR_fs8_ssee_z2_sigma}$ & $\val{bossR_fs8_ssee_z3}\pm\val{bossR_fs8_ssee_z3_sigma}$ \\
$\Lambda$CDM               & $\val{bossR_fs8_lcdm_z1}\pm\val{bossR_fs8_lcdm_z1_sigma}$ & $\val{bossR_fs8_lcdm_z2}\pm\val{bossR_fs8_lcdm_z2_sigma}$ & $\val{bossR_fs8_lcdm_z3}\pm\val{bossR_fs8_lcdm_z3_sigma}$ \\
\citet{Alam2017}           & $0.497\pm0.045$   & $0.458\pm0.038$   & $0.436\pm0.034$ \\
\midrule
\SSEE{} vs.\ published     & $\val{bossR_pub_ssee_z1}\sigma$ & $\val{bossR_pub_ssee_z2}\sigma$ & $\val{bossR_pub_ssee_z3}\sigma$ \\
$\Lambda$CDM vs.\ published& $\val{bossR_pub_lcdm_z1}\sigma$ & $\val{bossR_pub_lcdm_z2}\sigma$ & $\val{bossR_pub_lcdm_z3}\sigma$ \\
\bottomrule
\end{tabular}
\end{table}

The fit quality is the headline:
\begin{center}
\begin{tabular}{lcc}
\toprule
 & $\ln(10^{10}A_s)$ & $\chi^2$ \\
\midrule
$\Lambda$CDM & $\val{bossR_logA_lcdm}\pm\val{bossR_logA_lcdm_sigma}$ & $\val{bossR_chi2_lcdm}$ \\
\SSEE{}      & $\val{bossR_logA_ssee}\pm\val{bossR_logA_ssee_sigma}$ & $\val{bossR_chi2_ssee}$ \\
\bottomrule
\end{tabular}
\end{center}
$\Delta\chi^2_{\SSEE-\Lambda{\rm CDM}} = \val{bossR_dchi2}$ over $222$ points
with the same $37$ free parameters, i.e.\ a mild preference for the algebraic
background, not a decisive one. Each model carries its own neutrino mass
($\sum m_\nu=\val{SUM_MNU_EV_d5}$\,eV in \SSEE{}, $0.06$\,eV in $\Lambda$CDM); an
earlier version of these chains ran \SSEE{} with the $\Lambda$CDM value and
found $\Delta\chi^2=+0.06$, so the correction
% ORIGEN-VALOR: +0.06 = 207.638-207.576, archive/logs_superados/R1R2_boss_lpt_cobaya_20260908_mnu006.json
moved the comparison towards
\SSEE{}, not away from it. The algebraic background describes raw galaxy
clustering at least as well as Planck's fitted one.

\paragraph{Method control.} Our own $\Lambda$CDM row is the control: if the
pipeline is sound, it must reproduce what BOSS published. It does, within
$\val{bossR_pub_lcdm_z1}\sigma$, $\val{bossR_pub_lcdm_z2}\sigma$ and $\val{bossR_pub_lcdm_z3}\sigma$ at the three redshifts. We flag
two caveats rather than let the reader infer agreement we have not earned.
First, $z=0.38$ at $\val{bossR_pub_lcdm_z1}\sigma$ is the weakest of the three and we do not
claim it as clean. Second, our uncertainties ($\pm0.020$) are less than half
the published ones ($\pm0.034$ to $\pm0.045$), because we hold the background
\emph{fixed} whereas \citet{Alam2017} marginalise over it: the smaller error
bar buys its precision by not asking a question, and the two are therefore
not the same quantity. A $\Lambda$CDM run with the background free --- the
version that would make the error bars comparable --- is the outstanding item
for this section.

% FUENTES (2026-10-01): D y r por \val (mide_As_o_producto, etapa DVC); perfil,
%   marginal y su distancia por \val (analiza_boss_R1R2, que lee boss_aisla_neutrinos.json).
%   chi2 BOSS con A_s del CMB = 198.07: results/logs/base_sin_particula.log, linea 75
%   (rejilla con el m_nu de SSEE, conjunta_tres_sondas.py)
\paragraph{Clustering does not measure the amplitude.} The marginal means
above, $\val{bossR_logA_ssee}\pm\val{bossR_logA_ssee_sigma}$ (\SSEE{}) and $\val{bossR_logA_lcdm}\pm\val{bossR_logA_lcdm_sigma}$ ($\Lambda$CDM), are
not a measurement of $A_s$, and we do not compare them with the CMB. Three
measurements show it. First, the marginal and conditional widths of
$\log(10^{10}A_s)$ differ by a factor $D=\sigma_{\rm marg}/\sigma_{\rm cond}=\val{bossR_D_ssee}$,
the amplitude being degenerate with the linear bias ($r=\val{bossR_r_ssee}$); for KiDS on the
same fixed background $D=\val{kidsR_D_ssee}$. Second, the central value moves with the
estimator: the profile minimum of the same data, same $k_{\max}$ and same
model sits at $\val{bossR_perfil_logA}\pm\val{bossR_perfil_logA_sigma}$, $\val{bossR_perfil_vs_marg}\sigma$ from the marginal mean --- a
volume effect of marginalising over the nuisance parameters, which biases the
amplitude low. Third, fixing the amplitude to the CMB value costs the
clustering $\chi^2$ only $0.63$ ($\val{bossR_perfil_chi2}\to198.07$ in the profile). Galaxy
clustering is therefore consistent with the CMB amplitude and does not vote on
it.

\paragraph{The displacement is not $\omega_c$ in disguise.}
\label{par:wc-profile}
Amplitude and cold dark matter density are degenerate in clustering: both
raise the small-scale power. A referee is entitled to ask whether the low
$\ln(10^{10}A_s)$ above is really a low $\omega_c$ wearing another name. We
answer with the symmetric experiment. Holding the amplitude \emph{fixed} at
the CMB value and profiling $\omega_c$ instead --- everything else in the
background still fixed, only the galaxy biases fitted --- BOSS asks for
$\omega_c = \val{pwc_ssee_cmb}\pm\val{pwc_ssee_cmb_sigma}$, which is $\val{pwc_ssee_cmb_dist}\sigma$ from the algebraic
$\mathrm{KAL_0}\,\omega_b\,n_s = \val{OMEGA_C_H2_d6}$. Returning the amplitude BOSS
itself prefers, $\omega_c$ comes back to $\val{pwc_ssee_boss}\pm\val{pwc_ssee_boss_sigma}$, i.e.\
$\val{pwc_ssee_boss_dist}\sigma$ from the algebraic value. Each point of these
profiles minimises the galaxy biases from five starting points and is swept
along the grid in both directions; ten independent random starts at the
minimum and at both ends of the grid fail to lower $\chi^2$ by more than
$0.01$. The pre-registered control is that
second profile: if $\omega_c$ had not returned, the first would have been
measuring the edge of the parametrisation rather than the data. It returns.

The same profile run on the $\Lambda$CDM background, with identical rigidity,
displaces in the same direction and by a comparable amount
($\val{pwc_lcdm_boss}\rightarrow\val{pwc_lcdm_cmb}$, i.e.\ $\val{pwc_lcdm_cmb_dist}\sigma$ from
Planck's $\omega_c=\val{pwc_lcdm_ref}$ once the CMB amplitude is imposed), so
the displacement is a property of the data and not manufactured by either
background. Two consequences follow. First, the low clustering amplitude is
not a disguised matter density. Second --- and this was not sought --- the
algebraic $\omega_c$ sits \emph{closer} to what a structure survey at
$z\simeq0.5$ asks for ($\val{pwc_ssee_boss_dist}\sigma$) than Planck's fitted $\Lambda$CDM value
does ($\val{pwc_lcdm_boss_dist}\sigma$).

We deliberately do not quote a $\chi^2$ comparison between models with
$\omega_c$ free. In \SSEE{} that parameter is fixed by algebra: releasing it
takes the model outside itself, so such a number would not describe \SSEE{}
and cannot arbitrate between the two frameworks. The profile is a diagnostic
of where the data pulls, not a fit. What it does license is the statement
that $\Lambda$CDM gains $\val{gan_wc_lcdm}$ in $\chi^2$ from that (equally illegitimate)
freedom against $\val{gan_wc_ssee}$ for \SSEE{} --- Planck's $\omega_c$ needs the
larger correction. The legitimate model comparison remains the one in the table
above, each framework run as itself.

\subsection{Why the exploratory linear-Kaiser stage does not publish}
\label{sec:limitation-kaiser}

An earlier, exploratory pass over the same data used linear Kaiser
redshift-space distortions rather than the one-loop model reported above.%
% 2026-09-07: la frase decia "The multipole model used for the results in
% this section IS linear Kaiser", en presente, y el ultimo parrafo de la
% misma subseccion dice que la corrida de produccion usa velocileptors. Las
% dos no podian ser ciertas a la vez. La segunda es la correcta. Measured against the published BOSS DR12
consensus values (Alam et al.~2017 \citep{Alam2017}) with the growth rate left free, as in the
official analysis, the pipeline reproduces the published values within
0.92$\sigma$ on average ($\chi^2$/dof$=1.048$; \S\ref{sec:boss-method}),
confirming the machinery --- data reading, window convolution, and
Alcock--Paczynski deformation --- is sound. However, the linear model
under-predicts $f\sigma_8$ by 9--15\% (increasing with $z$) once matched
against small-scale modes ($k>0.08\,h/{\rm Mpc}$), a known limitation of
first-order perturbation theory in the mildly non-linear regime. The
production run for this paper uses a one-loop
effective-theory model (velocileptors, Chen, Vlah \& White 2020 \citep{ChenVlahWhite2020}) that
resolves this limitation, and \S\ref{sec:fsigma8} above reports its output,
not the linear-theory exploratory result.

% ============================================================================
\section{Conclusions}
% ============================================================================

The growth sector of \SSEE{} has a single matter density,
$\Omega_m=\val{OMEGA_M_TOTAL_d6}$, entering both the background and the growth of structure
with no partition, no additional particle, and no non-minimal coupling
parameter. Measured against raw survey data --- rather than against
$\Lambda$CDM-conditioned summary statistics --- it has no $S_8$ tension. On
KiDS-Legacy the amplitude the shear asks for lies $0.49\sigma$ from the one
the model's own CMB fit supplies; with that amplitude fixed the growth sector
has no free cosmological parameter, the fit costs nothing
($\chi^2_{\rm min}=417.97$ against $418.34$), and $S_8=\val{S8_leg_pred}$ is predicted
against $\val{S8_legacy_dato}\pm\val{S8_legacy_dato_sigma}$ measured. On the earlier KiDS-1000 release the
algebraic background matches $\Lambda$CDM run through the same pipeline with
four more free parameters ($\Delta\chi^2=\val{dchi2_k1000}$, $\Delta\mathrm{BIC}=\val{dBIC_k1000}$).
Against raw BOSS DR12 clustering the algebraic background fits marginally
better ($\Delta\chi^2=\val{bossR_dchi2}$ at equal freedom, not a decisive
difference), and clustering is consistent with the CMB amplitude
without constraining it. The primordial amplitude is therefore counted once, at
the CMB, and carried unchanged to every late-time probe. Cosmic shear and
clustering do not discriminate between the backgrounds; what separates them is
the parameter count.

The methodological point (\S\ref{sec:method}) is offered as a general caution
independent of \SSEE{}: any model without background-level free parameters
cannot exploit the mechanism by which combined datasets average away their own
internal tensions, and is therefore a comparatively clean probe of whether
those tensions are real.

% ============================================================================
\clearpage
\bibliographystyle{plainnat}
\bibliography{ssee_paper6_growth}

\end{document}
